\documentclass[11pt]{article}

\usepackage[margin=1in]{geometry}
\usepackage{amsmath,amssymb,amsthm}
\usepackage{enumitem}
\usepackage{xcolor}
\usepackage{hyperref}

\hypersetup{
  colorlinks=true,
  linkcolor=blue!60!black,
  urlcolor=blue!60!black
}

\newcommand{\R}{\mathbb{R}}
\newcommand{\N}{\mathbb{N}}
\newcommand{\Q}{\mathbb{Q}}
\newcommand{\B}{\mathcal{B}}
\newcommand{\F}{\mathcal{F}}
\newcommand{\C}{\mathcal{C}}
\newcommand{\Prob}{\mathbb{P}}
\newcommand{\E}{\mathbb{E}}
\newcommand{\Var}{\operatorname{Var}}
\newcommand{\Cov}{\operatorname{Cov}}
\newcommand{\Knowledge}[1]{%
  \par\smallskip\noindent
  \textbf{Knowledge used.}
  {\small\color{blue!60!black}\raggedright #1\par}
  \par\smallskip
}

\title{Chapter 7 Exercises}
\author{\textsc{Solutions to Rick Durrett's}\\
\textit{Probability: Theory and Examples}}
\date{\today}

\begin{document}
\maketitle

\noindent
These solutions use the local Chapter 7 LLM/Viki knowledge set in
\path{Probability/LLM_Viki_Dataset/Chapter7_Knowledge}.  The cited
knowledge IDs refer to entries in
\path{chapter7_knowledge_pieces.jsonl}.

\section*{Exercise 7.1.1}

\noindent\textbf{Question.}
Given $s<t$, find
\[
  \Prob(B(s)>0,B(t)>0).
\]
Hint: write the probability in terms of $B(s)$ and $B(t)-B(s)$.

\Knowledge{
\path{durrett_ch7_brownian_definition};
\path{durrett_ch7_gaussian_process_characterization}.
}

\noindent\textbf{Solution.}
Write
\[
  X=B(s),\qquad Y=B(t)-B(s).
\]
Then $X$ and $Y$ are independent normal random variables with
$X\sim N(0,s)$ and $Y\sim N(0,t-s)$.  Let
\[
  U=\frac{X}{\sqrt{s}},\qquad V=\frac{Y}{\sqrt{t-s}}.
\]
Then $(U,V)$ has a rotationally symmetric standard normal distribution
on $\R^2$, and the desired event is
\[
  \left\{U>0,\ \sqrt{s}\,U+\sqrt{t-s}\,V>0\right\}.
\]
Because the joint density is radial, this probability is the angle of
the corresponding sector divided by $2\pi$.

The boundary lines are $U=0$ and
\[
  \sqrt{s}\,U+\sqrt{t-s}\,V=0,
  \qquad\text{or}\qquad
  V=-\sqrt{\frac{s}{t-s}}\,U.
\]
The sector satisfying both inequalities has angle
\[
  \frac{\pi}{2}+\arctan\sqrt{\frac{s}{t-s}}
  =
  \frac{\pi}{2}+\arcsin\sqrt{\frac{s}{t}}.
\]
Therefore
\[
  \boxed{
  \Prob(B(s)>0,B(t)>0)
  =
  \frac14+\frac{1}{2\pi}\arcsin\sqrt{\frac{s}{t}}
  }.
\]

\section*{Exercise 7.1.2}

\noindent\textbf{Question.}
Find
\[
  \E(B_1^2B_2B_3).
\]

\Knowledge{
\path{durrett_ch7_brownian_definition};
\path{durrett_ch7_gaussian_process_characterization}.
}

\noindent\textbf{Solution.}
Let
\[
  X=B_1,\qquad Y=B_2-B_1,\qquad Z=B_3-B_2.
\]
By independent increments, $X,Y,Z$ are independent $N(0,1)$ random
variables.  Then
\[
  B_2=X+Y,\qquad B_3=X+Y+Z,
\]
so
\[
  \E(B_1^2B_2B_3)
  =
  \E\left[X^2(X+Y)(X+Y+Z)\right].
\]
The terms involving the independent centered variable $Z$ have mean
zero, hence
\[
  \E(B_1^2B_2B_3)
  =
  \E\left[X^2(X+Y)^2\right].
\]
Expanding and using independence and centering,
\[
  \E\left[X^2(X+Y)^2\right]
  =
  \E X^4+2\E X^3\,\E Y+\E X^2\,\E Y^2
  =
  3+0+1
  =
  4.
\]
Thus
\[
  \boxed{\E(B_1^2B_2B_3)=4}.
\]

\section*{Exercise 7.1.3}

\noindent\textbf{Question.}
Let
\[
  W=\int_0^t B_s\,ds.
\]
Find $\E W$ and $\E W^2$.  What is the distribution of $W$?

\Knowledge{
\path{durrett_ch7_brownian_definition};
\path{durrett_ch7_gaussian_process_characterization}.
}

\noindent\textbf{Solution.}
Since $\E B_s=0$ for all $s$,
\[
  \E W
  =
  \int_0^t \E B_s\,ds
  =
  0.
\]
Also, using $\E(B_rB_s)=r\wedge s$,
\[
  \E W^2
  =
  \E\left(\int_0^t B_r\,dr\int_0^t B_s\,ds\right)
  =
  \int_0^t\int_0^t (r\wedge s)\,dr\,ds.
\]
By symmetry,
\[
  \int_0^t\int_0^t (r\wedge s)\,dr\,ds
  =
  2\int_0^t\int_0^s r\,dr\,ds
  =
  2\int_0^t \frac{s^2}{2}\,ds
  =
  \frac{t^3}{3}.
\]
Hence
\[
  \boxed{\E W=0,\qquad \E W^2=\frac{t^3}{3}}.
\]

Finally, $W$ is a limit in $L^2$ of Riemann sums
\[
  \sum_i B_{s_i}\Delta s_i,
\]
and each such sum is a centered Gaussian random variable because
Brownian motion is a Gaussian process.  Therefore its $L^2$ limit $W$ is
also Gaussian.  Thus
\[
  \boxed{W\sim N\left(0,\frac{t^3}{3}\right)}.
\]

\section*{Exercise 7.1.4}

\noindent\textbf{Question.}
Show that $A\in\F^o$ if and only if there is a sequence of times
$t_1,t_2,\ldots$ in $[0,\infty)$ and a Borel set
$B\in\B(\R^{\{1,2,\ldots\}})$ such that
\[
  A=\{\omega:(\omega(t_1),\omega(t_2),\ldots)\in B\}.
\]
In words, all events in $\F^o$ depend on only countably many
coordinates.

\Knowledge{
\path{durrett_ch7_kolmogorov_extension_construction};
\path{durrett_ch7_brownian_definition}.
}

\noindent\textbf{Solution.}
Let $\Omega^o$ be the space of all real-valued functions on
$[0,\infty)$, and let $\F^o$ be the $\sigma$-field generated by the
coordinate maps $\omega\mapsto\omega(t)$.

First suppose that
\[
  A=\{\omega:(\omega(t_1),\omega(t_2),\ldots)\in B\}
\]
for some sequence $t_1,t_2,\ldots$ and some
$B\in\B(\R^{\{1,2,\ldots\}})$.  The map
\[
  \omega\mapsto(\omega(t_1),\omega(t_2),\ldots)
\]
is $\F^o/\B(\R^{\{1,2,\ldots\}})$-measurable, since each coordinate is
$\F^o$-measurable.  Therefore $A\in\F^o$.

Conversely, let $\mathcal{G}$ be the collection of all subsets of
$\Omega^o$ that depend on only countably many coordinates in the above
sense.  We claim that $\mathcal{G}$ is a $\sigma$-field.  Complements
are immediate: if $A$ is represented using $(t_n)$ and $B$, then $A^c$
is represented using the same coordinates and $B^c$.

Now let $A_1,A_2,\ldots\in\mathcal{G}$.  For each $j$, choose
coordinates $t_{j,1},t_{j,2},\ldots$ and a Borel set $B_j$ representing
$A_j$.  The union of the countable coordinate sets
\[
  \{t_{j,k}:j,k\ge 1\}
\]
is still countable, so enumerate it as $u_1,u_2,\ldots$.  Each event
$A_j$ can be rewritten as a Borel condition on
$(\omega(u_1),\omega(u_2),\ldots)$.  Hence
$\bigcup_j A_j$ is also a Borel condition on the same countable list of
coordinates.  Thus $\mathcal{G}$ is a $\sigma$-field.

Every finite-dimensional cylinder event belongs to $\mathcal{G}$, since
it depends on finitely many, hence countably many, coordinates.  Because
$\F^o$ is generated by the finite-dimensional cylinder events, we get
\[
  \F^o\subseteq\mathcal{G}.
\]
The first paragraph proved $\mathcal{G}\subseteq\F^o$, so
$\mathcal{G}=\F^o$.  This proves the claim.

\section*{Exercise 7.1.5}

\noindent\textbf{Question.}
Looking at the proof of Theorem 7.1.6 carefully shows that if
$\gamma>5/6$, then $B_t$ is not Holder continuous with exponent
$\gamma$ at any point in $[0,1]$.  Show, by considering $k$ increments
instead of $3$, that the last conclusion is true for all
\[
  \gamma>\frac12+\frac1k.
\]

\Knowledge{
\path{durrett_ch7_brownian_nowhere_differentiable};
\path{durrett_ch7_brownian_holder_half};
\path{durrett_ch7_brownian_definition}.
}

\noindent\textbf{Solution.}
Fix $C<\infty$ and an integer $k\ge 3$.  For $n\ge k$, define the event
$A_n(C)$ that there is an $s\in[0,1]$ such that
\[
  |B_t-B_s|\le C|t-s|^\gamma
  \qquad\text{whenever } |t-s|\le k/n.
\]
If Brownian motion is Holder continuous of exponent $\gamma$ at some
point with local constant $C$, then $A_n(C)$ occurs for all large $n$.

Now divide $[0,1]$ into intervals of length $1/n$.  If $A_n(C)$ occurs,
then for some block of $k$ consecutive increments near $s$, each
increment must be small.  More precisely, for a constant $C_k$ depending
only on $C$ and $k$, the event $A_n(C)$ is contained in
\[
  \bigcup_{j=0}^{n-k}
  \bigcap_{\ell=1}^k
  \left\{
  \left|B_{(j+\ell)/n}-B_{(j+\ell-1)/n}\right|
  \le C_k n^{-\gamma}
  \right\}.
\]
The $k$ increments in a fixed block are independent $N(0,1/n)$ random
variables.  If $Z\sim N(0,1/n)$, then
\[
  \Prob(|Z|\le C_k n^{-\gamma})
  =
  \Prob(|N(0,1)|\le C_k n^{1/2-\gamma})
  \le C_k' n^{1/2-\gamma}
\]
for another constant $C_k'$.  Therefore
\[
  \Prob(A_n(C))
  \le
  n\left(C_k' n^{1/2-\gamma}\right)^k
  =
  (C_k')^k n^{1+k(1/2-\gamma)}.
\]
If
\[
  \gamma>\frac12+\frac1k,
\]
then $1+k(1/2-\gamma)<0$, so $\Prob(A_n(C))\to 0$.

Thus, for each fixed $C$, the probability that such a Holder point with
constant $C$ exists is $0$.  Taking the countable union over
$C=1,2,\ldots$ gives probability $0$ for the existence of any Holder
point of exponent $\gamma$ in $[0,1]$.  Hence, almost surely, Brownian
motion is not Holder continuous of exponent $\gamma$ at any point of
$[0,1]$ whenever $\gamma>1/2+1/k$.

\section*{Exercise 7.1.6}

\noindent\textbf{Question.}
Fix $t$ and let
\[
  \Delta_{m,n}
  =
  B(tm2^{-n})-B(t(m-1)2^{-n}).
\]
Compute
\[
  \E\left(\sum_{m\le 2^n}\Delta_{m,n}^2-t\right)^2
\]
and use Borel--Cantelli to conclude that
\[
  \sum_{m\le 2^n}\Delta_{m,n}^2\to t
  \qquad\text{a.s. as }n\to\infty.
\]

\Knowledge{
\path{durrett_ch7_quadratic_variation_hint};
\path{durrett_ch7_brownian_definition}.
}

\noindent\textbf{Solution.}
For fixed $n$, the increments
\[
  \Delta_{1,n},\ldots,\Delta_{2^n,n}
\]
are independent and each has distribution $N(0,t2^{-n})$.  Hence
\[
  \E\Delta_{m,n}^2=t2^{-n},
  \qquad
  \E\Delta_{m,n}^4=3(t2^{-n})^2.
\]
Therefore
\[
  \Var(\Delta_{m,n}^2)
  =
  \E\Delta_{m,n}^4-(\E\Delta_{m,n}^2)^2
  =
  2t^2 2^{-2n}.
\]
Since the squared increments are independent,
\[
  \E\left(\sum_{m\le 2^n}\Delta_{m,n}^2\right)
  =
  2^n t2^{-n}
  =
  t,
\]
and
\[
  \E\left(\sum_{m\le 2^n}\Delta_{m,n}^2-t\right)^2
  =
  \Var\left(\sum_{m\le 2^n}\Delta_{m,n}^2\right)
  =
  2^n\cdot 2t^2 2^{-2n}
  =
  \boxed{2t^2 2^{-n}}.
\]

Now let
\[
  Q_n=\sum_{m\le 2^n}\Delta_{m,n}^2.
\]
For any $\varepsilon>0$, Chebyshev's inequality gives
\[
  \Prob(|Q_n-t|>\varepsilon)
  \le
  \frac{\E(Q_n-t)^2}{\varepsilon^2}
  =
  \frac{2t^2}{\varepsilon^2}2^{-n}.
\]
Since
\[
  \sum_{n=1}^\infty \Prob(|Q_n-t|>\varepsilon)<\infty,
\]
the Borel--Cantelli lemma implies that
$|Q_n-t|>\varepsilon$ occurs only finitely often, almost surely.  Taking
$\varepsilon$ over the positive rationals gives
\[
  \boxed{
  \sum_{m\le 2^n}\Delta_{m,n}^2\to t
  \quad\text{almost surely}.
  }
\]

\section*{Exercise 7.2.1}

\noindent\textbf{Question.}
Let
\[
  T_0=\inf\{s>0:B_s=0\}
  \qquad\text{and}\qquad
  R=\inf\{t>1:B_t=0\}.
\]
$R$ is for right or return.  Use the Markov property at time $1$ to get
\[
  \Prob_x(R>1+t)
  =
  \int p_1(x,y)\Prob_y(T_0>t)\,dy.
\]

\Knowledge{
\path{durrett_ch7_markov_property};
\path{durrett_ch7_immediate_positive_negative_hits};
\path{durrett_ch7_one_dimensional_recurrence}.
}

\noindent\textbf{Solution.}
The event $\{R>1+t\}$ says that, after time $1$, Brownian motion does
not hit $0$ during the next time interval of length $t$.  If we condition
on $B_1=y$, then the Markov property says that the post-$1$ process has
the law of Brownian motion started at $y$.  Therefore
\[
  \Prob_x(R>1+t\mid B_1)
  =
  \Prob_{B_1}(T_0>t).
\]
Taking expectations gives
\[
  \Prob_x(R>1+t)
  =
  \E_x\Prob_{B_1}(T_0>t).
\]
Since $B_1$ has density $p_1(x,y)$ under $\Prob_x$, this becomes
\[
  \boxed{
  \Prob_x(R>1+t)
  =
  \int_{\R} p_1(x,y)\Prob_y(T_0>t)\,dy
  }.
\]
When $y=0$, $\Prob_0(T_0>t)=0$ for $t>0$ by immediate return, which is
consistent with the displayed formula.

\section*{Exercise 7.2.2}

\noindent\textbf{Question.}
Let
\[
  T_0=\inf\{s>0:B_s=0\}
  \qquad\text{and}\qquad
  L=\sup\{t\le 1:B_t=0\}.
\]
$L$ is for left or last.  Use the Markov property at time $0<t<1$ to
conclude
\[
  \Prob_0(L\le t)
  =
  \int p_t(0,y)\Prob_y(T_0>1-t)\,dy.
\]

\Knowledge{
\path{durrett_ch7_markov_property};
\path{durrett_ch7_immediate_positive_negative_hits}.
}

\noindent\textbf{Solution.}
The event $\{L\le t\}$ means that there is no zero of the Brownian path
in the time interval $(t,1]$.  Conditional on $B_t=y$, the future
process
\[
  B_{t+u},\qquad 0\le u\le 1-t,
\]
has, by the Markov property, the law of Brownian motion started from
$y$.  Thus
\[
  \Prob_0(L\le t\mid B_t)
  =
  \Prob_{B_t}(T_0>1-t).
\]
Taking expectations and using the transition density of $B_t$ under
$\Prob_0$ gives
\[
  \boxed{
  \Prob_0(L\le t)
  =
  \int_{\R} p_t(0,y)\Prob_y(T_0>1-t)\,dy
  }.
\]
Again the value at $y=0$ causes no difficulty, since
$\Prob_0(T_0>1-t)=0$.

\section*{Exercise 7.2.3}

\noindent\textbf{Question.}
If $a<b$, then with probability one $a$ is the limit of local maxima of
$B_t$ in $(a,b)$.  So the set of local maxima of $B_t$ is almost surely
dense.  However, unlike the zero set, it is countable.

\Knowledge{
\path{durrett_ch7_markov_property};
\path{durrett_ch7_blumenthal_zero_one};
\path{durrett_ch7_immediate_positive_negative_hits};
\path{durrett_ch7_brownian_definition}.
}

\noindent\textbf{Solution.}
Fix $a<b$.  Choose deterministic intervals
\[
  [u_n,v_n]\subset(a,b),
  \qquad
  u_n\downarrow a,\quad v_n\downarrow a,\quad u_n<v_n.
\]
For instance, for large $n$ one may take
$u_n=a+(b-a)/(n+1)$ and $v_n=a+(b-a)/n$.
By continuity, Brownian motion attains its maximum on each compact
interval $[u_n,v_n]$; let $\tau_n$ be a time at which this maximum is
attained.

We claim that, with probability one, $\tau_n$ is not an endpoint.  At
the left endpoint $u_n$, the increment process
\[
  B_{u_n+s}-B_{u_n},\qquad s\ge 0,
\]
is Brownian motion started from $0$.  By immediate crossing,
it is positive at arbitrarily small positive times almost surely.  Hence
$u_n$ cannot be the maximum point of $[u_n,v_n]$.

At the right endpoint $v_n$, consider the reversed increment process
\[
  X_s=B_{v_n-s}-B_{v_n},
  \qquad 0\le s\le v_n-u_n.
\]
Its finite-dimensional distributions are those of Brownian motion
started from $0$ on this finite interval.  Therefore it is also positive
at arbitrarily small positive times almost surely.  Equivalently,
there are times just to the left of $v_n$ at which
$B$ is larger than $B_{v_n}$.  Thus $v_n$ cannot be the maximum point
of $[u_n,v_n]$.

Therefore $\tau_n\in(u_n,v_n)$ almost surely.  Since $B_{\tau_n}$ is the
maximum on the whole interval $[u_n,v_n]$, the time $\tau_n$ is a local
maximum of the Brownian path.  Also $\tau_n\to a$, because
$u_n,v_n\to a$.  This proves that $a$ is almost surely a limit of local
maxima lying in $(a,b)$.

Applying the result to all rational intervals shows that local maxima
are almost surely dense in $[0,\infty)$.  The countability comment is
consistent with this: for Brownian paths, local maxima are almost surely
strict, and strict local maxima of a continuous function form a
countable set by assigning to each one a rational interval on which it
is the unique maximum.

\section*{Exercise 7.2.4}

\noindent\textbf{Question.}
\begin{enumerate}[label=(\roman*)]
\item Suppose $f(t)>0$ for all $t>0$.  Use Theorem 7.2.3 to conclude
that
\[
  \limsup_{t\downarrow 0}\frac{B(t)}{f(t)}=c
  \qquad \Prob_0\text{-a.s.},
\]
where $c\in[0,\infty]$ is a constant.
\item Show that if $f(t)=\sqrt{t}$, then $c=\infty$, so with probability
one Brownian paths are not Holder continuous of order $1/2$ at $0$.
\end{enumerate}

\Knowledge{
\path{durrett_ch7_blumenthal_zero_one};
\path{durrett_ch7_immediate_positive_negative_hits};
\path{durrett_ch7_brownian_scaling};
\path{durrett_ch7_brownian_holder_half}.
}

\noindent\textbf{Solution.}
Let
\[
  L=\limsup_{t\downarrow 0}\frac{B(t)}{f(t)}.
\]
For each real $r$, the event $\{L\le r\}$ depends only on the values of
Brownian motion in arbitrarily small intervals after time $0$.  Hence
\[
  \{L\le r\}\in\F_0^+.
\]
By Blumenthal's $0$-$1$ law,
\[
  \Prob_0(L\le r)\in\{0,1\}
  \qquad\text{for every }r.
\]
Since Brownian motion started at $0$ immediately visits positive values,
$L\ge 0$ almost surely.  Define
\[
  c=\inf\{r\in\Q:r\ge 0,\ \Prob_0(L\le r)=1\},
\]
with the convention $c=\infty$ if the set is empty.  The preceding
$0$-$1$ statement implies that $L=c$ almost surely.  This proves (i).

Now take $f(t)=\sqrt{t}$.  Fix $M>0$ and $\delta>0$.  Since
$B_\delta/\sqrt{\delta}\sim N(0,1)$,
\[
  \Prob_0\left(
    \sup_{0<u\le \delta}\frac{B_u}{\sqrt{u}}>M
  \right)
  \ge
  \Prob_0\left(\frac{B_\delta}{\sqrt{\delta}}>M\right)
  =
  \Prob(N(0,1)>M)>0.
\]
Letting $\delta\downarrow 0$ through a decreasing sequence gives
\[
  \Prob_0\left(
    \limsup_{t\downarrow 0}\frac{B_t}{\sqrt{t}}\ge M
  \right)>0.
\]
This event lies in $\F_0^+$, so Blumenthal's $0$-$1$ law upgrades its
probability to $1$.  Since $M$ is arbitrary,
\[
  \limsup_{t\downarrow 0}\frac{B_t}{\sqrt{t}}=\infty
  \qquad\Prob_0\text{-a.s.}
\]
Thus $c=\infty$.

If the Brownian path were Holder continuous of order $1/2$ at $0$, then
there would be constants $C,\eta>0$ such that
\[
  |B_t-B_0|\le C\sqrt{t}
  \qquad\text{for }0<t<\eta.
\]
Since $B_0=0$, this would make $B_t/\sqrt{t}$ bounded above near $0$,
contradicting the infinite limsup just proved.  Therefore Brownian
paths are almost surely not Holder continuous of order $1/2$ at $0$.

\section*{Exercise 7.3.1}

\noindent\textbf{Question.}
Let $A$ be an $F_\sigma$, that is, a countable union of closed sets.
Show that
\[
  T_A=\inf\{t:B_t\in A\}
\]
is a stopping time.

\Knowledge{
\path{durrett_ch7_stopping_time_definition};
\path{durrett_ch7_hitting_times_open_closed};
\path{durrett_ch7_stopping_time_limits}.
}

\noindent\textbf{Solution.}
Write
\[
  A=\bigcup_{n=1}^\infty K_n,
\]
where each $K_n$ is closed.  Let
\[
  T_n=\inf\{t:B_t\in K_n\}.
\]
By the closed-set hitting-time theorem, each $T_n$ is a stopping time.
Moreover,
\[
  T_A=\inf_{n\ge 1}T_n,
\]
because hitting the union is the same as hitting at least one of the
closed sets.

For every $t\ge 0$,
\[
  \{T_A<t\}
  =
  \left\{\inf_n T_n<t\right\}
  =
  \bigcup_{n=1}^\infty \{T_n<t\}.
\]
Each event $\{T_n<t\}$ is in $\F_t$, and $\F_t$ is closed under
countable unions.  Hence $\{T_A<t\}\in\F_t$ for all $t$, so $T_A$ is a
stopping time.

\section*{Exercise 7.3.2}

\noindent\textbf{Question.}
If $S$ and $T$ are stopping times, then
\[
  S\wedge T=\min\{S,T\},\qquad
  S\vee T=\max\{S,T\},\qquad
  S+T
\]
are also stopping times.  In particular, if $t\ge 0$, then
$S\wedge t$, $S\vee t$, and $S+t$ are stopping times.

\Knowledge{
\path{durrett_ch7_stopping_time_definition};
\path{durrett_ch7_stopping_time_limits}.
}

\noindent\textbf{Solution.}
For the minimum,
\[
  \{S\wedge T<t\}
  =
  \{S<t\}\cup\{T<t\}\in\F_t.
\]
Thus $S\wedge T$ is a stopping time.

For the maximum, use the equivalent criterion
$\{U\le t\}\in\F_t$ for stopping times under a right-continuous
filtration.  Then
\[
  \{S\vee T\le t\}
  =
  \{S\le t\}\cap\{T\le t\}\in\F_t.
\]
So $S\vee T$ is a stopping time.

For the sum, for $t>0$,
\[
  \{S+T<t\}
  =
  \bigcup_{\substack{q\in\Q\\0<q<t}}
  \{S<q\}\cap\{T<t-q\}.
\]
Indeed, if $S+T<t$, choose a rational $q$ with
$S<q<t-T$.  Conversely, any such $q$ implies $S+T<t$.  The event
$\{S<q\}$ lies in $\F_q\subseteq\F_t$, and
$\{T<t-q\}$ lies in $\F_{t-q}\subseteq\F_t$, so the countable union is
in $\F_t$.  Therefore $S+T$ is a stopping time.

Taking $T$ to be the deterministic stopping time identically equal to
$t$ gives the conclusions for $S\wedge t$, $S\vee t$, and $S+t$.

\section*{Exercise 7.3.3}

\noindent\textbf{Question.}
Let $T_n$ be a sequence of stopping times.  Show that
\[
  \sup_n T_n,\qquad
  \inf_n T_n,\qquad
  \limsup_n T_n,\qquad
  \liminf_n T_n
\]
are stopping times.

\Knowledge{
\path{durrett_ch7_stopping_time_definition};
\path{durrett_ch7_stopping_time_limits}.
}

\noindent\textbf{Solution.}
For the infimum,
\[
  \left\{\inf_n T_n<t\right\}
  =
  \bigcup_{n=1}^\infty \{T_n<t\}\in\F_t,
\]
so $\inf_n T_n$ is a stopping time.

For the supremum, use the equivalent $\le$ formulation:
\[
  \left\{\sup_n T_n\le t\right\}
  =
  \bigcap_{n=1}^\infty \{T_n\le t\}\in\F_t.
\]
Thus $\sup_n T_n$ is a stopping time.

Now write
\[
  \liminf_{n\to\infty}T_n
  =
  \sup_{m\ge 1}\inf_{n\ge m}T_n,
  \qquad
  \limsup_{n\to\infty}T_n
  =
  \inf_{m\ge 1}\sup_{n\ge m}T_n.
\]
The inner infima and suprema are stopping times by the first two parts,
and applying the same two closure properties once more shows that
$\liminf_n T_n$ and $\limsup_n T_n$ are stopping times.

\section*{Exercise 7.3.4}

\noindent\textbf{Question.}
Let $S$ be a stopping time, let $A\in\F_S$, and let
\[
  R=
  \begin{cases}
  S, & \text{on }A,\\
  \infty, & \text{on }A^c.
  \end{cases}
\]
Show that $R$ is a stopping time.

\Knowledge{
\path{durrett_ch7_stopping_time_definition};
\path{durrett_ch7_random_shift_and_F_S}.
}

\noindent\textbf{Solution.}
For each $t\ge 0$,
\[
  \{R\le t\}
  =
  A\cap\{S\le t\}.
\]
Since $A\in\F_S$, the defining property of $\F_S$ gives
\[
  A\cap\{S\le t\}\in\F_t.
\]
Therefore $\{R\le t\}\in\F_t$ for every $t$, and hence $R$ is a
stopping time.

\section*{Exercise 7.3.5}

\noindent\textbf{Question.}
Let $S$ and $T$ be stopping times.  Show that
\[
  \{S<T\},\qquad \{S>T\},\qquad \{S=T\}
\]
are in $\F_S$ and in $\F_T$.

\Knowledge{
\path{durrett_ch7_stopping_time_definition};
\path{durrett_ch7_random_shift_and_F_S}.
}

\noindent\textbf{Solution.}
We use the equivalent characterization of $\F_S$ with
$A\cap\{S<t\}\in\F_t$ for every $t$.

First,
\[
  \{S<T\}\cap\{S<t\}
  =
  \bigcup_{\substack{q\in\Q\\q<t}}
  \{S<q\}\cap\{q<T\}.
\]
The event $\{S<q\}$ is in $\F_q$, and
$\{q<T\}=\{T\le q\}^c$ is also in $\F_q$.  Since $q<t$, each term lies
in $\F_t$, so the countable union lies in $\F_t$.  Hence
$\{S<T\}\in\F_S$.

Next,
\[
  \{S>T\}\cap\{S<t\}
  =
  \bigcup_{\substack{q\in\Q\\q<t}}
  \{T<q\}\cap\{q<S\}\cap\{S<t\}.
\]
Each event in the intersection is measurable by time $t$, so
$\{S>T\}\in\F_S$.  Finally,
\[
  \{S=T\}
  =
  \{S<T\}^c\cap\{S>T\}^c,
\]
so $\{S=T\}\in\F_S$ as well.

The same argument with the roles of $S$ and $T$ interchanged shows that
all three events also belong to $\F_T$.

\section*{Exercise 7.3.6}

\noindent\textbf{Question.}
Let $S$ and $T$ be stopping times.  Show that
\[
  \F_S\cap\F_T=\F_{S\wedge T}.
\]

\Knowledge{
\path{durrett_ch7_stopping_time_definition};
\path{durrett_ch7_random_shift_and_F_S};
\path{durrett_ch7_stopping_time_limits}.
}

\noindent\textbf{Solution.}
Since $S\wedge T\le S$ and $S\wedge T\le T$, information available at
time $S\wedge T$ is available at both later stopping times.  Thus
\[
  \F_{S\wedge T}\subseteq \F_S\cap\F_T.
\]

For the reverse inclusion, let $A\in\F_S\cap\F_T$.  We must show that
$A\in\F_{S\wedge T}$.  For every $t\ge 0$,
\[
  A\cap\{S\wedge T\le t\}
  =
  A\cap(\{S\le t\}\cup\{T\le t\}).
\]
Hence
\[
  A\cap\{S\wedge T\le t\}
  =
  (A\cap\{S\le t\})\cup(A\cap\{T\le t\}).
\]
Because $A\in\F_S$, the first term is in $\F_t$; because $A\in\F_T$,
the second term is in $\F_t$.  Therefore their union is in $\F_t$.
This proves $A\in\F_{S\wedge T}$.

Combining the two inclusions gives
\[
  \boxed{\F_S\cap\F_T=\F_{S\wedge T}}.
\]

\section*{Exercise 7.4.1}

\noindent\textbf{Question.}
Let $B_t=(B_t^1,B_t^2)$ be a two-dimensional Brownian motion starting
from $0$.  Let
\[
  T_a=\inf\{t:B_t^2=a\}.
\]
Show that $B^1(T_a)$, $a\ge 0$, has independent increments and
\[
  B^1(T_a)\stackrel{d}{=}aB^1(T_1).
\]
Use this to conclude that $B^1(T_a)$ has a Cauchy distribution.

\Knowledge{
\path{durrett_ch7_d_dimensional_brownian};
\path{durrett_ch7_hitting_time_subordinator};
\path{durrett_ch7_hitting_time_scaling};
\path{durrett_ch7_hitting_time_density}.
}

\noindent\textbf{Solution.}
The two coordinates $B^1$ and $B^2$ are independent one-dimensional
Brownian motions.  The random time $T_a$ is determined only by the
second coordinate.

Let
\[
  0=a_0<a_1<\cdots<a_n.
\]
Conditioning on the hitting times
$T_{a_0},T_{a_1},\ldots,T_{a_n}$, the increments
\[
  B^1(T_{a_j})-B^1(T_{a_{j-1}}),
  \qquad 1\le j\le n,
\]
are independent centered normal random variables with variances
$T_{a_j}-T_{a_{j-1}}$.  The process $\{T_a:a\ge 0\}$ has stationary
independent increments as a function of the level $a$, so averaging over
the hitting times preserves independence.  Therefore
\[
  X_a:=B^1(T_a),\qquad a\ge 0,
\]
has independent increments.

Next compute its characteristic function.  Since $B^1$ is independent
of $T_a$,
\[
  \E e^{iuB^1(T_a)}
  =
  \E\left[\E\left(e^{iuB^1(T_a)}\mid T_a\right)\right]
  =
  \E e^{-u^2T_a/2}.
\]
Brownian hitting-time scaling gives $T_a\stackrel{d}{=}a^2T_1$, hence
\[
  \E e^{iuB^1(T_a)}
  =
  \E e^{-(au)^2T_1/2}
  =
  \E e^{iu\,aB^1(T_1)}.
\]
Thus
\[
  \boxed{B^1(T_a)\stackrel{d}{=}aB^1(T_1)}.
\]

Finally, using the first-passage density
\[
  \Prob(T_a\in ds)
  =
  \frac{a}{\sqrt{2\pi s^3}}e^{-a^2/(2s)}\,ds,
  \qquad a>0,
\]
or equivalently its Laplace transform, we get
\[
  \E e^{-u^2T_a/2}=e^{-a|u|}.
\]
This is the characteristic function of the centered Cauchy distribution
with scale parameter $a$, whose density is
\[
  \boxed{
  f_a(x)=\frac{1}{\pi}\frac{a}{a^2+x^2},
  \qquad x\in\R,\ a>0.
  }
\]
For $a=0$, $B^1(T_0)=0$ is degenerate.

\section*{Exercise 7.4.2}

\noindent\textbf{Question.}
Use (7.2.3) to show that
\[
  R=\inf\{t>1:B_t=0\}
\]
has probability density
\[
  \Prob_0(R\in d(1+t))/dt
  =
  \frac{1}{\pi t^{1/2}(1+t)}.
\]

\Knowledge{
\path{durrett_ch7_markov_property};
\path{durrett_ch7_hitting_time_density};
\path{durrett_ch7_immediate_positive_negative_hits}.
}

\noindent\textbf{Solution.}
From Exercise 7.2.1 with $x=0$,
\[
  \Prob_0(R>1+t)
  =
  \int_{\R}p_1(0,y)\Prob_y(T_0>t)\,dy.
\]
For $y\ne 0$, the hitting time of $0$ by Brownian motion started at
$y$ has the same law as the hitting time of $|y|$ by Brownian motion
started at $0$.  Hence its density is
\[
  \frac{|y|}{\sqrt{2\pi r^3}}e^{-y^2/(2r)},\qquad r>0.
\]
Therefore the density of $R$ at $1+t$ is obtained by differentiating the
survival function:
\[
  f_R(1+t)
  =
  \int_{\R}p_1(0,y)
  \frac{|y|}{\sqrt{2\pi t^3}}e^{-y^2/(2t)}\,dy.
\]
Substitute $p_1(0,y)=(2\pi)^{-1/2}e^{-y^2/2}$:
\[
  f_R(1+t)
  =
  \frac{1}{2\pi t^{3/2}}
  \int_{\R}|y|
  \exp\left\{-\frac{1+t}{2t}y^2\right\}\,dy.
\]
Since
\[
  \int_{\R}|y|e^{-cy^2}\,dy=\frac{1}{c},
  \qquad c>0,
\]
with $c=(1+t)/(2t)$, we obtain
\[
  f_R(1+t)
  =
  \frac{1}{2\pi t^{3/2}}\cdot\frac{2t}{1+t}
  =
  \boxed{\frac{1}{\pi t^{1/2}(1+t)}}.
\]

\section*{Exercise 7.4.3}

\noindent\textbf{Question.}
\begin{enumerate}[label=(\alph*)]
\item Generalize the proof of (7.4.5) to conclude that if $u<v\le a$,
then
\[
  \Prob_0(T_a<t,\ u<B_t<v)
  =
  \Prob_0(2a-v<B_t<2a-u).
\]
\item Let $(u,v)$ shrink down to $x$ in the preceding identity to
conclude that, for $x<a$,
\[
  \Prob_0(T_a<t,\ B_t\in dx)
  =
  p_t(0,2a-x)\,dx.
\]
\item Use $\{T_a<t\}=\{M_t>a\}$, where
$M_t=\max_{0\le s\le t}B_s$, and differentiate with respect to $a$ to
obtain the joint density
\[
  f_{(M_t,B_t)}(a,x)
  =
  \frac{2(2a-x)}{\sqrt{2\pi t^3}}
  e^{-(2a-x)^2/(2t)}.
\]
\end{enumerate}

\Knowledge{
\path{durrett_ch7_reflection_principle};
\path{durrett_ch7_strong_markov_property};
\path{durrett_ch7_hitting_time_density}.
}

\noindent\textbf{Solution.}
For part (a), condition at the first hitting time $T_a$.  On
$\{T_a<t\}$, the post-$T_a$ process
\[
  B_{T_a+r}-a,\qquad r\ge 0,
\]
is Brownian motion started from $0$, independent of the past.  By
symmetry, after the first hit of $a$, paths ending below $a$ in the
interval $(u,v)$ can be reflected across the horizontal line $a$ and
matched with paths whose endpoint lies in $(2a-v,2a-u)$.  Therefore
\[
  \boxed{
  \Prob_0(T_a<t,\ u<B_t<v)
  =
  \Prob_0(2a-v<B_t<2a-u)
  }.
\]

For part (b), divide by $v-u$ and let $u,v\to x<a$.  Since $B_t$ has
density $p_t(0,\cdot)$,
\[
  \boxed{
  \Prob_0(T_a<t,\ B_t\in dx)
  =
  p_t(0,2a-x)\,dx,
  \qquad x<a.
  }
\]
This is the subprobability density of $B_t$ on the event that the path
has crossed $a$ before time $t$ and ends below $a$.

For part (c), note that $\{T_a<t\}=\{M_t>a\}$.  Thus for $x<a$,
\[
  \Prob_0(M_t>a,\ B_t\in dx)
  =
  p_t(0,2a-x)\,dx.
\]
The joint density of $(M_t,B_t)$ is the negative derivative in $a$:
\[
  f_{(M_t,B_t)}(a,x)
  =
  -\frac{\partial}{\partial a}p_t(0,2a-x).
\]
Since
\[
  p_t(0,z)=\frac{1}{\sqrt{2\pi t}}e^{-z^2/(2t)},
\]
we have
\[
  -\frac{\partial}{\partial a}p_t(0,2a-x)
  =
  \frac{2(2a-x)}{t}p_t(0,2a-x).
\]
Therefore
\[
  \boxed{
  f_{(M_t,B_t)}(a,x)
  =
  \frac{2(2a-x)}{\sqrt{2\pi t^3}}
  e^{-(2a-x)^2/(2t)}
  },
\]
for $a\ge 0$ and $x<a$.

\section*{Exercise 7.4.4}

\noindent\textbf{Question.}
Let $A_{s,t}$ be the event that Brownian motion has at least one zero in
$[s,t]$.  Show that
\[
  \Prob_0(A_{s,t})
  =
  \frac{2}{\pi}\arccos\sqrt{\frac{s}{t}}.
\]
Hint: To get started note that
\[
  \Prob_0(A_{s,t})
  =
  2\int_0^\infty (2\pi s)^{-1/2}e^{-x^2/(2s)}
    \Prob_x(T_0\le t-s)\,dx
\]
and use the formula for the density function of $T_x$.  At the end,
you may need some trigonometry to convert arctangent into arccosine.

\Knowledge{
\path{durrett_ch7_markov_property};
\path{durrett_ch7_hitting_time_density};
\path{durrett_ch7_zero_set_structure}.
}

\noindent\textbf{Solution.}
Condition on $B_s=x$.  By symmetry and the Markov property,
\[
  \Prob_0(A_{s,t})
  =
  2\int_0^\infty (2\pi s)^{-1/2}e^{-x^2/(2s)}
    \Prob_x(T_0\le t-s)\,dx.
\]
For $x>0$, the density of $T_0$ under $\Prob_x$ is the same as the
density of $T_x$ under $\Prob_0$:
\[
  \Prob_x(T_0\in dr)
  =
  \frac{x}{\sqrt{2\pi r^3}}e^{-x^2/(2r)}\,dr.
\]
Thus
\[
  \Prob_0(A_{s,t})
  =
  2\int_0^\infty (2\pi s)^{-1/2}e^{-x^2/(2s)}
  \int_0^{t-s}
  \frac{x}{\sqrt{2\pi r^3}}e^{-x^2/(2r)}\,dr\,dx.
\]
Interchanging the order of integration gives
\[
  \Prob_0(A_{s,t})
  =
  \frac{1}{\pi\sqrt{s}}
  \int_0^{t-s}r^{-3/2}
  \int_0^\infty
  x\exp\left\{-\frac{x^2}{2}\left(\frac1s+\frac1r\right)\right\}
  dx\,dr.
\]
The inner integral is
\[
  \frac{sr}{s+r},
\]
so
\[
  \Prob_0(A_{s,t})
  =
  \frac{\sqrt{s}}{\pi}
  \int_0^{t-s}\frac{r^{-1/2}}{s+r}\,dr.
\]
Let $u=\sqrt{r/s}$.  Then $r=su^2$ and $dr=2su\,du$, so
\[
  \frac{\sqrt{s}}{\pi}
  \int_0^{t-s}\frac{r^{-1/2}}{s+r}\,dr
  =
  \frac{2}{\pi}
  \int_0^{\sqrt{(t-s)/s}}\frac{du}{1+u^2}.
\]
Therefore
\[
  \Prob_0(A_{s,t})
  =
  \frac{2}{\pi}\arctan\sqrt{\frac{t-s}{s}}.
\]
If $\theta=\arctan\sqrt{(t-s)/s}$, then
\[
  \cos\theta=\frac{1}{\sqrt{1+(t-s)/s}}=\sqrt{\frac{s}{t}},
\]
so $\theta=\arccos\sqrt{s/t}$.  Hence
\[
  \boxed{
  \Prob_0(A_{s,t})
  =
  \frac{2}{\pi}\arccos\sqrt{\frac{s}{t}}
  }.
\]

\section*{Exercise 7.5.1}

\noindent\textbf{Question.}
Let
\[
  T=\inf\{t:B_t\notin(-a,a)\}.
\]
Show that
\[
  \E_0 e^{-\lambda T}
  =
  \frac{1}{\cosh(a\sqrt{2\lambda})}.
\]

\Knowledge{
\path{durrett_ch7_optional_stopping_bounded};
\path{durrett_ch7_exponential_martingale};
\path{durrett_ch7_expected_interval_exit_time}.
}

\noindent\textbf{Solution.}
Let
\[
  \alpha=\sqrt{2\lambda}.
\]
Since both
\[
  e^{\alpha B_t-\alpha^2t/2}
  \qquad\text{and}\qquad
  e^{-\alpha B_t-\alpha^2t/2}
\]
are martingales, their average
\[
  M_t=e^{-\lambda t}\cosh(\alpha B_t)
\]
is also a martingale.  Apply optional stopping to $T\wedge n$:
\[
  1=M_0=\E_0\left[e^{-\lambda(T\wedge n)}
      \cosh(\alpha B_{T\wedge n})\right].
\]
Letting $n\to\infty$, the bounded convergence argument on the stopped
exit event gives $B_T=\pm a$ and hence
\[
  1
  =
  \E_0\left[e^{-\lambda T}\cosh(\alpha B_T)\right]
  =
  \cosh(\alpha a)\E_0 e^{-\lambda T}.
\]
Therefore
\[
  \boxed{
  \E_0 e^{-\lambda T}
  =
  \frac{1}{\cosh(a\sqrt{2\lambda})}
  }.
\]

\section*{Exercise 7.5.2}

\noindent\textbf{Question.}
The point of this exercise is to get information about the amount of
time it takes Brownian motion with drift $-b$,
\[
  X_t=B_t-bt,
\]
to hit level $a$.  Let
\[
  \tau=\inf\{t:B_t=a+bt\},
  \qquad a>0.
\]
Use the martingale $e^{\theta B_t-\theta^2t/2}$ with
\[
  \theta=b+\sqrt{b^2+2\lambda}
\]
to show
\[
  \E_0 e^{-\lambda\tau}
  =
  \exp\{-a(b+\sqrt{b^2+2\lambda})\}.
\]
Letting $\lambda\to0$ gives
\[
  \Prob_0(\tau<\infty)=e^{-2ab}.
\]

\Knowledge{
\path{durrett_ch7_exponential_martingale};
\path{durrett_ch7_optional_stopping_bounded}.
}

\noindent\textbf{Solution.}
Let
\[
  \theta=b+\sqrt{b^2+2\lambda}.
\]
Then
\[
  \frac{\theta^2}{2}-b\theta=\lambda,
  \qquad\text{or equivalently}\qquad
  b\theta-\frac{\theta^2}{2}=-\lambda.
\]
The exponential martingale
\[
  M_t=e^{\theta B_t-\theta^2t/2}
\]
stopped at $\tau\wedge n$ satisfies
\[
  1=\E_0 M_{\tau\wedge n}.
\]
On $\{\tau\le n\}$, $B_\tau=a+b\tau$, so
\[
  M_\tau
  =
  \exp\{\theta a+(b\theta-\theta^2/2)\tau\}
  =
  e^{\theta a}e^{-\lambda\tau}.
\]
On $\{\tau>n\}$, we have $B_n<a+bn$, and therefore
\[
  M_n
  \le
  \exp\{\theta a+(b\theta-\theta^2/2)n\}
  =
  e^{\theta a}e^{-\lambda n}.
\]
Letting $n\to\infty$ gives
\[
  1=e^{\theta a}\E_0 e^{-\lambda\tau},
\]
and hence
\[
  \boxed{
  \E_0 e^{-\lambda\tau}
  =
  e^{-a(b+\sqrt{b^2+2\lambda})}
  }.
\]
For $b>0$, letting $\lambda\downarrow0$ gives
\[
  \boxed{
  \Prob_0(\tau<\infty)=e^{-2ab}
  }.
\]

\section*{Exercise 7.5.3}

\noindent\textbf{Question.}
Let
\[
  \sigma=\inf\{t:B_t\notin(a,b)\},
  \qquad \lambda>0.
\]
Use the strong Markov property to show
\[
  \E_x e^{-\lambda T_a}
  =
  \E_x(e^{-\lambda\sigma};T_a<T_b)
  +
  \E_x(e^{-\lambda\sigma};T_b<T_a)\E_b e^{-\lambda T_a}.
\]
Interchange the roles of $a$ and $b$ to get a second equation, use
Theorem 7.5.7, and solve to get
\[
  \E_x(e^{-\lambda\sigma};T_a<T_b)
  =
  \frac{\sinh(\sqrt{2\lambda}(b-x))}
       {\sinh(\sqrt{2\lambda}(b-a))}
\]
and
\[
  \E_x(e^{-\lambda\sigma};T_b<T_a)
  =
  \frac{\sinh(\sqrt{2\lambda}(x-a))}
       {\sinh(\sqrt{2\lambda}(b-a))}.
\]

\Knowledge{
\path{durrett_ch7_strong_markov_property};
\path{durrett_ch7_hitting_time_laplace_exact}.
}

\noindent\textbf{Solution.}
Let
\[
  A_x=\E_x(e^{-\lambda\sigma};T_a<T_b),
  \qquad
  B_x=\E_x(e^{-\lambda\sigma};T_b<T_a),
\]
and put $\alpha=\sqrt{2\lambda}$.  On $\{T_a<T_b\}$, we have
$T_a=\sigma$.  On $\{T_b<T_a\}$, the strong Markov property at $\sigma$
gives
\[
  T_a=\sigma+T_a\circ\theta_\sigma,
\]
where the restarted Brownian motion begins at $b$.  Hence
\[
  \E_x e^{-\lambda T_a}
  =
  A_x+B_x\E_b e^{-\lambda T_a}.
\]
Using the one-sided hitting transform,
\[
  \E_x e^{-\lambda T_a}=e^{-\alpha(x-a)},
  \qquad
  \E_b e^{-\lambda T_a}=e^{-\alpha(b-a)}.
\]
Thus
\[
  e^{-\alpha(x-a)}
  =
  A_x+e^{-\alpha(b-a)}B_x.
\]
Interchanging $a$ and $b$ similarly gives
\[
  e^{-\alpha(b-x)}
  =
  B_x+e^{-\alpha(b-a)}A_x.
\]
Solving the two linear equations yields
\[
  A_x
  =
  \frac{e^{-\alpha(x-a)}
        -e^{-\alpha(b-a)}e^{-\alpha(b-x)}}
       {1-e^{-2\alpha(b-a)}}.
\]
Multiplying numerator and denominator by $e^{\alpha(b-a)}$ gives
\[
  A_x
  =
  \frac{e^{\alpha(b-x)}-e^{-\alpha(b-x)}}
       {e^{\alpha(b-a)}-e^{-\alpha(b-a)}}
  =
  \boxed{
  \frac{\sinh(\alpha(b-x))}{\sinh(\alpha(b-a))}
  }.
\]
The same algebra gives
\[
  \boxed{
  B_x
  =
  \frac{\sinh(\alpha(x-a))}{\sinh(\alpha(b-a))}
  }.
\]

\section*{Exercise 7.5.4}

\noindent\textbf{Question.}
If
\[
  T=\inf\{t:B_t\notin(a,b)\},
  \qquad a<0<b,\quad a\ne -b,
\]
then $T$ and $B_T^2$ are not independent, so we cannot calculate
$\E T^2$ as in the proof of Theorem 7.5.9.  Use the Cauchy--Schwarz
inequality to estimate $\E(TB_T^2)$ and conclude
\[
  \text{(i) }\E T^2\le 4\E B_T^4,
  \qquad
  \text{(ii) }\E B_T^4\le 36\E T^2.
\]

\Knowledge{
\path{durrett_ch7_heat_equation_martingales};
\path{durrett_ch7_optional_stopping_bounded};
\path{durrett_ch7_expected_interval_exit_time}.
}

\noindent\textbf{Solution.}
The polynomial
\[
  u(t,x)=x^4-6tx^2+3t^2
\]
satisfies $u_t+\frac12u_{xx}=0$, so
\[
  B_t^4-6tB_t^2+3t^2
\]
is a Brownian martingale.  Stopping at $T\wedge n$ and passing to the
limit gives
\[
  \E B_T^4-6\E(TB_T^2)+3\E T^2=0.
\]
Let
\[
  X=(\E T^2)^{1/2},
  \qquad
  Y=(\E B_T^4)^{1/2}.
\]
By Cauchy--Schwarz,
\[
  \E(TB_T^2)\le XY.
\]
From the martingale identity,
\[
  3X^2=6\E(TB_T^2)-Y^2\le 6XY.
\]
Hence $X\le 2Y$, and therefore
\[
  \boxed{\E T^2\le 4\E B_T^4}.
\]
The same identity also gives
\[
  Y^2=6\E(TB_T^2)-3X^2\le 6XY.
\]
Thus $Y\le 6X$, and
\[
  \boxed{\E B_T^4\le 36\E T^2}.
\]

\section*{Exercise 7.5.5}

\noindent\textbf{Question.}
Find a martingale of the form
\[
  B_t^6-c_1tB_t^4+c_2t^2B_t^2-c_3t^3
\]
and use it to compute the third moment of
\[
  T=\inf\{t:B_t\notin(-a,a)\}.
\]

\Knowledge{
\path{durrett_ch7_heat_equation_martingales};
\path{durrett_ch7_symmetric_exit_second_moment};
\path{durrett_ch7_optional_stopping_bounded}.
}

\noindent\textbf{Solution.}
Let
\[
  u(t,x)=x^6-c_1tx^4+c_2t^2x^2-c_3t^3.
\]
We choose the constants so that $u_t+\frac12u_{xx}=0$.  Now
\[
  u_t=-c_1x^4+2c_2tx^2-3c_3t^2,
\]
and
\[
  \frac12u_{xx}=15x^4-6c_1tx^2+c_2t^2.
\]
Equating coefficients gives
\[
  c_1=15,\qquad c_2=45,\qquad c_3=15.
\]
Thus
\[
  \boxed{
  B_t^6-15tB_t^4+45t^2B_t^2-15t^3
  }
\]
is a martingale.

For $T=\inf\{t:B_t\notin(-a,a)\}$, we have $B_T=\pm a$.  Optional
stopping and passage to the limit give
\[
  0
  =
  \E\left[B_T^6-15TB_T^4+45T^2B_T^2-15T^3\right].
\]
Therefore
\[
  0
  =
  a^6-15a^4\E T+45a^2\E T^2-15\E T^3.
\]
Using $\E T=a^2$ and $\E T^2=5a^4/3$,
\[
  0
  =
  a^6-15a^6+75a^6-15\E T^3
  =
  61a^6-15\E T^3.
\]
Hence
\[
  \boxed{\E T^3=\frac{61}{15}a^6}.
\]

\section*{Exercise 7.5.6}

\noindent\textbf{Question.}
Show that
\[
  (1+t)^{-1/2}\exp\left(\frac{B_t^2}{2(1+t)}\right)
\]
is a martingale and use this to conclude that
\[
  \limsup_{t\to\infty}
  \frac{B_t}{\{(1+t)\log(1+t)\}^{1/2}}
  \le 1
  \qquad\text{a.s.}
\]

\Knowledge{
\path{durrett_ch7_heat_equation_martingales};
\path{durrett_ch7_optional_stopping_bounded};
\path{durrett_ch7_brownian_martingale}.
}

\noindent\textbf{Solution.}
Let
\[
  u(t,x)=(1+t)^{-1/2}\exp\left(\frac{x^2}{2(1+t)}\right).
\]
A direct calculation gives
\[
  u_x=\frac{x}{1+t}u,\qquad
  u_{xx}=\left(\frac{1}{1+t}+\frac{x^2}{(1+t)^2}\right)u,
\]
and
\[
  u_t=\left(-\frac{1}{2(1+t)}
     -\frac{x^2}{2(1+t)^2}\right)u.
\]
Hence
\[
  u_t+\frac12u_{xx}=0.
\]
By the heat-equation martingale criterion,
\[
  M_t=(1+t)^{-1/2}\exp\left(\frac{B_t^2}{2(1+t)}\right)
\]
is a positive martingale with $M_0=1$.

Fix $\varepsilon>0$ and define
\[
  c_\varepsilon=\frac{(1+\varepsilon)^2-1}{2}>0.
\]
If for some $t\in[e^n,e^{n+1}]$,
\[
  B_t>(1+\varepsilon)\{(1+t)\log(1+t)\}^{1/2},
\]
then
\[
  M_t
  \ge
  (1+t)^{c_\varepsilon}
  \ge
  (1+e^n)^{c_\varepsilon}.
\]
Doob's maximal inequality for the nonnegative martingale $M_t$ gives
\[
  \Prob\left(
    \sup_{0\le s\le e^{n+1}}M_s\ge (1+e^n)^{c_\varepsilon}
  \right)
  \le
  (1+e^n)^{-c_\varepsilon}.
\]
The right side is summable in $n$.  By Borel--Cantelli, with
probability one, only finitely many intervals $[e^n,e^{n+1}]$ contain a
time $t$ for which
\[
  B_t>(1+\varepsilon)\{(1+t)\log(1+t)\}^{1/2}.
\]
Therefore
\[
  \limsup_{t\to\infty}
  \frac{B_t}{\{(1+t)\log(1+t)\}^{1/2}}
  \le 1+\varepsilon
  \qquad\text{a.s.}
\]
Letting $\varepsilon\downarrow0$ through rational values yields
\[
  \boxed{
  \limsup_{t\to\infty}
  \frac{B_t}{\{(1+t)\log(1+t)\}^{1/2}}
  \le 1
  \qquad\text{a.s.}
  }
\]

\section*{Exercise 7.6.1}

\noindent\textbf{Question.}
Suppose $f\in C^2$.  Show that $f(B_t)$ is a martingale if and only if
\[
  f(x)=a+bx
\]
for constants $a,b$.

\Knowledge{
\path{durrett_ch7_ito_formula_one_dimensional};
\path{durrett_ch7_stochastic_integral_is_martingale};
\path{durrett_ch7_heat_equation_martingales}.
}

\noindent\textbf{Solution.}
If $f(x)=a+bx$, then
\[
  f(B_t)=a+bB_t,
\]
which is a martingale because Brownian motion is a martingale.

Conversely, suppose $f(B_t)$ is a martingale.  Ito's formula gives
\[
  f(B_t)
  =
  f(B_0)+\int_0^t f'(B_s)\,dB_s
  +\frac12\int_0^t f''(B_s)\,ds.
\]
The stochastic integral is a local martingale.  Hence the finite
variation term
\[
  \frac12\int_0^t f''(B_s)\,ds
\]
must also be a local martingale.  A continuous finite variation local
martingale is constant, so
\[
  \int_0^t f''(B_s)\,ds=0
  \qquad\text{for all }t.
\]
Taking expectations under $\Prob_x$ and using the Brownian transition
density,
\[
  0
  =
  \int_0^t\int_{\R} f''(y)p_s(x,y)\,dy\,ds.
\]
Since $f''$ is continuous and the heat kernel is strictly positive, this
forces $f''\equiv0$.  Therefore $f$ is affine:
\[
  \boxed{f(x)=a+bx}.
\]

\section*{Exercise 7.6.2}

\noindent\textbf{Question.}
Show that
\[
  B_t^3-3tB_t
  \qquad\text{and}\qquad
  B_t^3-\int_0^t 3B_s\,ds
\]
are martingales.

\Knowledge{
\path{durrett_ch7_ito_formula_one_dimensional};
\path{durrett_ch7_ito_formula_time_space};
\path{durrett_ch7_stochastic_integral_is_martingale}.
}

\noindent\textbf{Solution.}
For the first process, let
\[
  u(t,x)=x^3-3tx.
\]
Then
\[
  u_t=-3x,\qquad u_x=3x^2-3t,\qquad u_{xx}=6x,
\]
so
\[
  u_t+\frac12u_{xx}=-3x+3x=0.
\]
By the time-space Ito formula,
\[
  u(t,B_t)=B_t^3-3tB_t
\]
is a martingale.

For the second process, apply one-dimensional Ito's formula to
$f(x)=x^3$:
\[
  B_t^3-B_0^3
  =
  \int_0^t 3B_s^2\,dB_s+\frac12\int_0^t 6B_s\,ds.
\]
Since $B_0=0$,
\[
  B_t^3-\int_0^t3B_s\,ds
  =
  \int_0^t3B_s^2\,dB_s,
\]
which is a martingale by the stochastic-integral martingale criterion.

\section*{Exercise 7.6.3}

\noindent\textbf{Question.}
Let
\[
  \beta_{2k}(t)=\E_0 B_t^{2k}.
\]
Use Ito's formula to relate $\beta_{2k}(t)$ to
$\beta_{2k-2}(t)$ and use this relationship to derive a formula for
$\beta_{2k}(t)$.

\Knowledge{
\path{durrett_ch7_ito_formula_one_dimensional};
\path{durrett_ch7_stochastic_integral_is_martingale}.
}

\noindent\textbf{Solution.}
Apply Ito's formula to $f(x)=x^{2k}$:
\[
  B_t^{2k}
  =
  \int_0^t 2kB_s^{2k-1}\,dB_s
  +
  \frac12\int_0^t 2k(2k-1)B_s^{2k-2}\,ds.
\]
Taking expectations gives
\[
  \beta_{2k}(t)
  =
  k(2k-1)\int_0^t \beta_{2k-2}(s)\,ds.
\]
Equivalently,
\[
  \beta_{2k}'(t)=k(2k-1)\beta_{2k-2}(t),
  \qquad
  \beta_0(t)=1.
\]
Induction now gives
\[
  \beta_{2k}(t)
  =
  (2k-1)(2k-3)\cdots 3\cdot 1\, t^k.
\]
Thus
\[
  \boxed{
  \beta_{2k}(t)
  =
  (2k-1)!!\,t^k
  =
  \frac{(2k)!}{2^k k!}\,t^k
  }.
\]

\section*{Exercise 7.6.4}

\noindent\textbf{Question.}
Apply Ito's formula to $|B_t|$.

\Knowledge{
\path{durrett_ch7_multidimensional_ito_formula};
\path{durrett_ch7_stochastic_integral_is_martingale}.
}

\noindent\textbf{Solution.}
For a $d$-dimensional Brownian motion, let
\[
  R_t=|B_t|.
\]
The function $x\mapsto |x|$ is not $C^2$ at $0$, so one applies Ito's
formula first to
\[
  f_\varepsilon(x)=(|x|^2+\varepsilon)^{1/2}
\]
and then lets $\varepsilon\downarrow0$.  Away from $0$,
\[
  \nabla |x|=\frac{x}{|x|},
  \qquad
  \Delta |x|=\frac{d-1}{|x|}.
\]
Thus, up to the first hitting time of $0$ and, by approximation, for
$d\ge2$ after starting from $0$,
\[
  \boxed{
  |B_t|
  =
  |B_0|
  +
  \sum_{i=1}^d\int_0^t \frac{B_s^i}{|B_s|}\,dB_s^i
  +
  \frac{d-1}{2}\int_0^t \frac{1}{|B_s|}\,ds
  }.
\]
The stochastic integral has quadratic variation $t$ whenever
$|B_s|>0$, since
\[
  \sum_{i=1}^d \left(\frac{B_s^i}{|B_s|}\right)^2=1.
\]
So the radial process has the familiar Bessel-form decomposition
\[
  d|B_t|=dW_t+\frac{d-1}{2|B_t|}\,dt,
\]
where $W_t$ is a one-dimensional Brownian motion generated by the radial
martingale part.  In dimension $d=1$, the singularity at $0$ contributes
the local-time term in Tanaka's formula.

\section*{Exercise 7.6.5}

\noindent\textbf{Question.}
Apply Ito's formula to $|B_t|^2$.  Use this to conclude that
\[
  \E_0|B_t|^2=td.
\]

\Knowledge{
\path{durrett_ch7_multidimensional_ito_formula};
\path{durrett_ch7_square_martingale}.
}

\noindent\textbf{Solution.}
Let $B_t=(B_t^1,\ldots,B_t^d)$ and
\[
  f(x)=|x|^2=\sum_{i=1}^d x_i^2.
\]
Then
\[
  \frac{\partial f}{\partial x_i}=2x_i,
  \qquad
  \frac{\partial^2 f}{\partial x_i^2}=2.
\]
The multidimensional Ito formula gives
\[
  |B_t|^2-|B_0|^2
  =
  \sum_{i=1}^d\int_0^t 2B_s^i\,dB_s^i
  +
  \frac12\int_0^t \sum_{i=1}^d 2\,ds.
\]
Since $B_0=0$,
\[
  |B_t|^2
  =
  2\sum_{i=1}^d\int_0^t B_s^i\,dB_s^i
  +
  dt.
\]
The stochastic integral has mean $0$, so
\[
  \boxed{\E_0|B_t|^2=dt}.
\]

\end{document}
